% !tp-example Basic usage
\documentclass{article}
\begin{document}
\wcDiv[data-tp-reference-scope]{
  A useful detail \rawhtml*|<tp-ref href="^detail"></tp-ref>|.

  \wcTpNote[ref=detail,title=Detail]{
    A note with \textbf{formatted content}.

  }
}
\end{document}

% !tp-example Notes, bibliography and glossary
\documentclass{article}
\begin{document}
\wcDiv[data-tp-reference-scope]{
  Read a note \rawhtml*|<tp-ref href="^method"></tp-ref>|, consult a source \wcTpRef*[href=@book]{}, or look up \rawhtml*|<tp-ref href="%ecosystem"></tp-ref>|. Compare another note \rawhtml*|<tp-ref href="^weather"></tp-ref>|, then read the first note again \rawhtml*|<tp-ref href="^method"></tp-ref>|.

  \wcTpNote[ref=method,title=Method]{
    Observe the habitat at \textbf{three different times} of day.

    \begin{itemize}
      \item Record the weather.
      \item Compare your observations.
    \end{itemize}
  }
  \wcTpBiblio[ref=book,title=Book]{
    Alex Example. \emph{Field observation handbook}. Example Press, 2026.

  }
  \wcTpGlossary[ref=ecosystem,title=Ecosystem]{
    A community of organisms interacting with their physical environment.

  }
  \wcTpNote[ref=weather]{
    Record the temperature and cloud cover.

  }
  \wcTpBiblio[ref=atlas]{
    Alex Example. \emph{A field atlas}. Example Press, 2025.

  }
  \wcTpGlossary[ref=community]{
    A group of interacting populations.

  }
  \subsection*{Notes}
  \wcTpListof[selector=tp-note]{}
  \subsection*{Bibliography}
  \wcTpListof[selector=tp-biblio]{}
  \subsection*{Glossary}
  \wcTpListof[selector=tp-glossary]{}
}
\end{document}
